% Section: Introduction
% Source: paper_draft_v0.1.md § 1. Introduction
% [VERIFY ACM TEMPLATE BEFORE SUBMISSION]

\section{Introduction}

Maintaining cyber hygiene---keeping privileged accounts active and monitored,
endpoint agents deployed, patches current, and telemetry fresh---is a persistent
operational challenge for security operations centers (SOCs), particularly in
public-sector environments subject to compliance mandates such as CISA Binding
Operational Directives~\cite{cisa_bod2301_2022} and NIST SP
800-40~\cite{nist_sp80040_2022}. Failures in hygiene posture create the enabling
conditions for identity-based attacks: dormant accounts reactivated by adversaries,
endpoints missing agent coverage, and stale privileged credentials that may be
invisible to behavioral analytics requiring recent baselines.

Despite the operational importance of hygiene-state monitoring, the anomaly-detection
research community lacks a dedicated public benchmark for this problem. Existing
benchmarks---the LANL Unified Host and Network Dataset~\cite{kent2015lanl}, CERT
Insider Threat datasets~\cite{glasser2013cert}, CICIDS~\cite{sharafaldin2018cicids},
and attack-emulation repositories such as Mordor/Security
Datasets~\cite{rodriguez2019mordor}---provide rich telemetry for network intrusion
or insider-threat detection but do not model identity hygiene state, patch posture,
vulnerability exposure, or telemetry freshness and missingness as first-class
evaluation axes. Commercial platforms (Microsoft Defender for Identity, Splunk UBA,
Exabeam) address related detection problems but are closed and do not release
benchmark datasets.

This gap matters for two reasons. First, researchers who develop anomaly-detection
methods for security operations have no way to compare their methods on
hygiene-specific telemetry under controlled conditions. Second, there is no
principled accounting of when unsupervised ML methods add value over simpler
rule-based approaches for hygiene anomalies---a question with real operational
implications, since rule maintenance is expensive and ML deployment requires
justification.

This paper makes the following contributions:

\begin{itemize}[leftmargin=*]

\item \textbf{C1: Synthetic generator.} A fully open, seeded synthetic telemetry
  generator covering eleven entity/event tables (Active Directory users, groups,
  computers, assets; login events; group membership events; endpoint patch state;
  vulnerability records; remediation events; account lifecycle events; telemetry
  freshness log) with documented structural priors.

\item \textbf{C2: Anomaly taxonomy.} Twelve cyber-hygiene anomaly classes
  (AH-01 through AH-12) with ATT\&CK enabling-condition mappings (not technique
  detection claims), injection protocols, and class imbalance specifications.

\item \textbf{C3: Benchmark task suite.} Seven evaluation tasks (T1--T7) under
  five telemetry conditions, with pre-registered $k$ values, feature sets, and
  method applicability flags.

\item \textbf{C4: Comparative evaluation, failure-aware.} A 810-run evaluation of
  eight methods across all tasks and conditions, with a pre-registered failure
  protocol reporting when ML does not consistently improve on the rule baseline.

\item \textbf{C5: Negative-result reporting.} Explicit reporting and analysis of
  the 86.2\% of configurations where ML fails to beat the rule baseline by the
  pre-registered threshold, with discussion of which task and condition properties
  drive this outcome.

\end{itemize}

The remainder of the paper is organized as follows. Section~\ref{sec:related}
reviews related benchmarks. Section~\ref{sec:design} describes the \hygienebench
design. Section~\ref{sec:setup} presents experimental setup.
Section~\ref{sec:results} reports results. Section~\ref{sec:discussion} discusses
findings and limitations. Section~\ref{sec:conclusion} concludes.
